% Evidence draft only; not submission-ready.
% This source requires XeLaTeX with xeCJK support. No LaTeX compiler was available locally;
% no compiled PDF from this source has been validated. The accepted reading PDF is
% artifacts/v2/reporting/rna_evidence_v2_manuscript_draft.pdf (SHA256=473eaa07…52d0b).
% Options for packages loaded elsewhere
\PassOptionsToPackage{unicode}{hyperref}
\PassOptionsToPackage{hyphens}{url}
%
\documentclass[
  11pt]{article}
\usepackage{amsmath,amssymb}
\usepackage{lmodern}
\usepackage{iftex}
\ifPDFTeX
  \usepackage[T1]{fontenc}
  \usepackage[utf8]{inputenc}
  \usepackage{textcomp} % provide euro and other symbols
\else % if luatex or xetex
  \usepackage{unicode-math}
  \usepackage{xeCJK}
  \defaultfontfeatures{Scale=MatchLowercase}
  \defaultfontfeatures[\rmfamily]{Ligatures=TeX,Scale=1}
\fi
% Use upquote if available, for straight quotes in verbatim environments
\IfFileExists{upquote.sty}{\usepackage{upquote}}{}
\IfFileExists{microtype.sty}{% use microtype if available
  \usepackage[]{microtype}
  \UseMicrotypeSet[protrusion]{basicmath} % disable protrusion for tt fonts
}{}
\makeatletter
\@ifundefined{KOMAClassName}{% if non-KOMA class
  \IfFileExists{parskip.sty}{%
    \usepackage{parskip}
  }{% else
    \setlength{\parindent}{0pt}
    \setlength{\parskip}{6pt plus 2pt minus 1pt}}
}{% if KOMA class
  \KOMAoptions{parskip=half}}
\makeatother
\usepackage{xcolor}
\IfFileExists{xurl.sty}{\usepackage{xurl}}{} % add URL line breaks if available
\IfFileExists{bookmark.sty}{\usepackage{bookmark}}{\usepackage{hyperref}}
\hypersetup{
  pdftitle={Contamination-aware evaluation of RNA sequence representations: pretraining exposure, homology, and limits of structure-assisted learning},
  pdfauthor={Anonymous evidence draft},
  hidelinks,
  pdfcreator={LaTeX via pandoc}}
\urlstyle{same} % disable monospaced font for URLs
\usepackage[margin=1in]{geometry}
\usepackage{longtable,booktabs,array}
\usepackage{calc} % for calculating minipage widths
% Correct order of tables after \paragraph or \subparagraph
\usepackage{etoolbox}
\makeatletter
\patchcmd\longtable{\par}{\if@noskipsec\mbox{}\fi\par}{}{}
\makeatother
% Allow footnotes in longtable head/foot
\IfFileExists{footnotehyper.sty}{\usepackage{footnotehyper}}{\usepackage{footnote}}
\makesavenoteenv{longtable}
\setlength{\emergencystretch}{3em} % prevent overfull lines
\providecommand{\tightlist}{%
  \setlength{\itemsep}{0pt}\setlength{\parskip}{0pt}}
\setcounter{secnumdepth}{-\maxdimen} % remove section numbering
\ifLuaTeX
  \usepackage{selnolig}  % disable illegal ligatures
\fi

\title{Contamination-aware evaluation of RNA sequence representations:
pretraining exposure, homology, and limits of structure-assisted
learning}
\author{Anonymous evidence draft}
\date{2026-11-18}

\begin{document}
\maketitle

\hypertarget{contamination-aware-evaluation-of-rna-sequence-representations-pretraining-exposure-homology-and-limits-of-structure-assisted-learning}{%
\section{Contamination-aware evaluation of RNA sequence representations:
pretraining exposure, homology, and limits of structure-assisted
learning}\label{contamination-aware-evaluation-of-rna-sequence-representations-pretraining-exposure-homology-and-limits-of-structure-assisted-learning}}

\textbf{Document status:} evidence draft, not submission-ready.
Numerical claims are tied to the artifacts and independent checks listed
in Section 7. No new architecture accuracy claim is made.

\hypertarget{abstract}{%
\subsection{Abstract}\label{abstract}}

Sequence-only evaluation of RNA representations can be difficult to
interpret when the pretraining corpus shares sequences with downstream
benchmarks. We therefore audited exact and approximate sequence exposure
before evaluating structure-assisted learning. Exact matching and
MMseqs2 searches at 80\% identity with at least 80\% bidirectional
coverage were performed for 165,854 unique queries against RNAcentral
22.0, Rfam 14.9, and Ensembl releases 109/56. At least one known source
yielded a qualifying match for 162,489 queries, including all 196 VL0
and all 5,388 bpRNA-new sequences. These results prohibit a clean
out-of-distribution interpretation, although they do not provide a
complete historical contamination denominator because the NCBI nt
snapshot is not pinned. Under this contamination-aware framework, the
tested dual-view JEPA did not outperform its structure-free host on
model-dev macro-F1 (0.844264 versus 0.851460). BPfold inner/outer motif
inputs also showed no gain over the corresponding host or shuffled
control. A frozen H0 representation achieved micro-AUPRC 0.7715 on
verified RNA--protein contacts, compared with 0.7216 for native RiNALMo
and 0.6292 for within-chain label shuffling. However, H2 shared H0's
frozen encoder, structure-assisted arms were memory-blocked, and short
chains remained weak. The results support explicit exposure reporting
and strict controls rather than a new accuracy or
independent-generalization claim.

\textbf{Keywords:} RNA secondary structure, pretraining contamination,
sequence homology, JEPA, RNA--protein interface, evaluation reliability

\hypertarget{1-introduction}{%
\subsection{1. Introduction}\label{1-introduction}}

RNA language models are commonly evaluated by withholding exact
sequences, sequence clusters, or Rfam families from downstream training.
Such splits address different notions of independence and cannot, by
themselves, establish that a pretrained representation has not
encountered a test sequence. Exact pretraining overlap, approximate
homology, and family identity must therefore be reported separately.

This work evaluates the evidentiary basis of three proposed improvements
to RNA representation learning: dual-view sequence--structure
prediction, BPfold-inspired motif inputs, and transfer to RNA--protein
interface prediction. The primary question was whether training-time
structural information could improve single-checkpoint sequence-only
inference. The analysis does not treat an ensemble as a new
architecture, canonical-nested pairs as three-dimensional structures, or
a downstream probe as a substitute for the primary accuracy test.

The contributions are: (i) a versioned, hash-bound audit of exact
sequence and approximate-homology exposure across three public source
families; (ii) an official-metadata family mapping that preserves
missing and ambiguous labels; (iii) a separation of engineering
correctness, model-dev evidence, and external conclusions for
structure-assisted and motif models; and (iv) independently checked
interface baselines with verified contact labels and explicit
pretraining-exposure limits. The outcome is an evidence framework and a
set of negative and qualified results, not a claim of improved accuracy.

\hypertarget{2-source-identity-and-evaluation-scope}{%
\subsection{2. Source identity and evaluation
scope}\label{2-source-identity-and-evaluation-scope}}

Table 1 identifies the public sources used for exposure analysis. These
sources were used only for auditing, not for additional model training.
RNAcentral and Rfam exact matches were computed from versioned FASTA
files. Ensembl 109 and Ensembl Bacteria 56 comprised 31,046 ncRNA files,
701,138,004 compressed bytes, and 4,040,439 records, of which 4,032,030
passed the original 16--8192 nt length filter. Independent checks
covered file identity, SHA256, gzip integrity, nonempty FASTA records,
and counts.

\textbf{Table 1. Public sources and roles}

\begin{longtable}[]{@{}lll@{}}
\toprule
Source & Version & Role \\
\midrule
\endhead
RNAcentral & 22.0 & Exact and approximate-homology exposure \\
Rfam & 14.9 FASTA, \texttt{full\_region}, \texttt{family} & Exposure and
official family mapping \\
Ensembl & 109 and Bacteria 56 & Exact and approximate-homology
exposure \\
NCBI nt & Historical snapshot not pinned & Unresolved source; no
current-file substitution \\
\bottomrule
\end{longtable}

RiNALMo preprocessing was rechecked against the paper: sequences shorter
than 16 or longer than 8192 nt were removed, and inputs longer than 1022
nt were randomly cropped during training. Thus, a full-length source
match does not establish that every nucleotide of a long sequence
appeared in every training crop. No such uncertainty is used to dismiss
direct matches.

The historical xens2 ensemble remains a baseline, not a new single-model
architecture. Mol-JEPA was used only as a design reference at commit
\texttt{27f0065b69c73d18eec7a32ecc413b6dc3f33d90}, under CC BY-NC 4.0.
The implementation did not add torch-geometric. The original ResNet2D,
Turner parameters, and exact noncrossing dynamic-programming decoder
were preserved in the relevant model comparisons.

\hypertarget{3-methods}{%
\subsection{3. Methods}\label{3-methods}}

\hypertarget{31-exact-exposure-and-query-identity}{%
\subsubsection{3.1 Exact exposure and query
identity}\label{31-exact-exposure-and-query-identity}}

Sequences were uppercased and T was normalized to U; ambiguous
nucleotide symbols were otherwise preserved. Exact overlap was computed
after applying the 16--8192 nt source-length filter. Duplicate rows were
distinguished from unique sequences. Query names were not assumed to be
unique identifiers: after finding same-name/different-sequence
conflicts, candidate sequences were checked directly against the
original FASTA sources. Each source's unique and row-level counts had to
agree with its initial audit before union counts were accepted.

The three exact-source unions were deduplicated by normalized full
sequence rather than added. Homology searches used MMseqs2 17.b804f
\texttt{easy-search}, \texttt{-\/-search-type\ 3},
\texttt{-\/-min-seq-id\ 0.8}, \texttt{-c\ 0.8},
\texttt{-\/-cov-mode\ 0}, and \texttt{-\/-max-seqs\ 300}. The latter was
not treated as a cap on final reported target rows. Every reported hit
was independently required to satisfy
\texttt{alnlen\ \textgreater{}=\ 0.8*qlen} and
\texttt{alnlen\ \textgreater{}=\ 0.8*tlen}. Nine RNAcentral alignments
failed this target-coverage check, but no query was lost from the
three-source union: 162,489 queries were detected both before and after
the postfilter. A query without an accepted match is marked ``not
detected,'' not ``nonhomologous.'' Searches were run with a 50 GiB
memory split. RNAcentral initially failed once with an estimated 185 GiB
prefilter requirement and again during competition with the Ensembl
search. It was subsequently run serially with the same threshold and
complete query set, preserving both failure directories.

\hypertarget{32-official-family-mapping}{%
\subsubsection{3.2 Official family
mapping}\label{32-official-family-mapping}}

Rfam families were propagated only through exact full-sequence matches
linking a query to a FASTA accession/coordinate region and then to
\texttt{full\_region.rfam\_acc} and \texttt{family.rfam\_acc}. Multiple
families for the same sequence were marked ambiguous rather than
resolved arbitrarily. FASTA regions lacking official family metadata
were reported separately from sequences absent from Rfam 14.9. Rfam 14.9
does not cover later releases, and TR1c includes mixed sources without
universal Rfam labels. Consequently, mapped-family overlap is a partial
analysis and is not interchangeable with sequence homology.

\hypertarget{33-structure-assisted-and-motif-evaluations}{%
\subsubsection{3.3 Structure-assisted and motif
evaluations}\label{33-structure-assisted-and-motif-evaluations}}

H0--H6 used the same initialization, data, and primary loss. H0 was the
no-JEPA host; H1 used dual-view MSE; H2 added target SIGReg; H3 shuffled
structure labels; H4 used a fixed random structure encoder; H5 was a
parameter-matched MLP; and H6 used a global mean structure vector. Pilot
selection used single-seed model-dev macro-F1, not the external test
sets. Independent pair counting checked micro/macro-F1, precision, and
recall. A separate native-initialization comparison was explicitly
distinguished from Plan-A warm-start and from optimizer continuation.

The BPfold experiment used the official MIT \texttt{key.energy} file and
its inner/outer definitions. It compared no motif, inner/outer motif
input, shuffled motif, Turner=0, and energy-only input. It was not
presented as a complete BPfold reproduction or as a Turner/RIPE
experiment.

\hypertarget{34-rnaprotein-contact-probes}{%
\subsubsection{3.4 RNA--protein contact
probes}\label{34-rnaprotein-contact-probes}}

Contacts were recomputed from original mmCIF polymer and atom records.
An observed RNA residue was labeled positive when a heavy atom of a
polypeptide chain was within 5 Å. Unobserved residues were not treated
as negatives. Complete polymer/reference-sequence correspondence and
modified-parent mappings were hash-checked. The final train/holdout sets
contained 804/202 chains and 1,453,871/231,155 observed residues, with
no PDB, 80\%-cluster, or exact-sequence overlap across the split. These
downstream exclusions do not remove pretraining exposure.

Six frozen-representation or composition controls were run with seeds
42, 2026, and 3407: native RiNALMo, H0, local handcrafted features,
within-chain label shuffling, length, and GC content. All 18 runs
stopped at an internal-validation plateau. Thresholds were chosen only
using PDB- and sequence-cluster-isolated training-internal validation.
Micro-AUPRC and binary macro-F1 were recomputed independently with
scikit-learn; distinct score thresholds made AUPRC independent of
tied-row ordering. H2 was excluded as an independent JEPA arm because
its encoder was tensor-identical to H0 at the selected checkpoint.
Structure-assisted arms remained blocked by the original ResNet2D's
full-pair memory requirement at 7,224 nt.

\hypertarget{4-results}{%
\subsection{4. Results}\label{4-results}}

\hypertarget{41-exact-and-approximate-homology-exposure}{%
\subsubsection{4.1 Exact and approximate-homology
exposure}\label{41-exact-and-approximate-homology-exposure}}

\textbf{Table 2. Exact overlap and detected 80\% homology after
within-source deduplication}

\begin{longtable}[]{@{}lrr@{}}
\toprule
Set & Exact union / total & Homology union / queries \\
\midrule
\endhead
TR1c & 37,870 / 42,564 & 40,283 / 42,564 \\
VL0 & 193 / 196 & 196 / 196 \\
bpRNA-new & 5,381 / 5,388 & 5,388 / 5,388 \\
TS0 & 1,253 / 1,288 & 1,281 / 1,288 \\
Rfam & 10,327 / 10,790 & 10,711 / 10,790 \\
ncRNA test & 24,912 / 25,038 & 25,012 / 25,038 \\
Interface holdout & 86 / 202 & 111 / 198 \\
\bottomrule
\end{longtable}

Across all queries, 162,489/165,854 had a qualifying match in at least
one known source; 3,365 did not. The 198 interface queries are the
length-filtered subset of 202 chains. The exact and homology columns
answer different questions and are not expected to be identical. No-hit
records cannot be treated as clean because of search heuristics, the
target cap, and the unresolved nt snapshot.

A read-only reanalysis of the verified original xens2 replay on
bpRNA-new gave macro-F1 0.542209. The exact-match-detected subset had
macro-F1 0.542424 (n=5,381), compared with 0.377253 for seven sequences
without an exact match. All 5,388 sequences had at least one qualifying
homology match, leaving no homology-not-detected comparison arm. The
seven-sequence subset is not a clean control and is insufficient for
significance testing or a causal exposure-effect claim. An older
alternative result with 385 different full records was not mixed with
the P0-verified historical/replay identity.

\hypertarget{42-family-mapping-is-not-equivalent-to-pretraining-independence}{%
\subsubsection{4.2 Family mapping is not equivalent to pretraining
independence}\label{42-family-mapping-is-not-equivalent-to-pretraining-independence}}

The independent mapping check covered 58,742 rows. Rfam evaluation
contained 8,002 uniquely mapped rows, 110 ambiguous rows, 2,111 rows
with FASTA matches but no family metadata, and 567 without an exact Rfam
14.9 match. Corresponding counts for bpRNA-new were 4,938, 0, 442, and
8, and for TR1c were 24,392, 35, 4,487, and 13,650. The uniquely mapped
family sets contained 1,597, 477, and 1,732 families, respectively.

Rfam evaluation and TR1c shared 120 mapped families. The bpRNA-new
intersection was zero only for the mapped downstream labels. This zero
supports the intended downstream family split; it does not negate the
exact and approximate pretraining exposure in Table 2.

\hypertarget{43-dual-view-pilot-and-initialization-confounding}{%
\subsubsection{4.3 Dual-view pilot and initialization
confounding}\label{43-dual-view-pilot-and-initialization-confounding}}

\textbf{Table 3. Model-dev macro-F1 for the dual-view pilot}

\begin{longtable}[]{@{}lr@{}}
\toprule
Arm & Macro-F1 \\
\midrule
\endhead
H0: no JEPA & 0.851460 \\
H1: dual-view MSE & 0.851186 \\
H2: MSE + target SIGReg & 0.844264 \\
H3: shuffled structure & 0.839411 \\
H4: fixed random structure encoder & 0.848010 \\
H5: parameter-matched MLP & 0.847129 \\
H6: global mean structure vector & 0.847233 \\
\bottomrule
\end{longtable}

H2 trailed H0 by 0.0071958679 and did not exceed the shuffle, random,
MLP, or global-mean controls. The preregistered decision was
\texttt{STOP\_PILOT\_NEGATIVE\_DEV\_ONLY}; three-seed formal training
and frozen external evaluation were not started. A small micro-F1
increase was not substituted for the primary macro-F1 decision.

Native-initialization H0/H2 values were 0.653980 and 0.658859,
respectively, a difference of +0.004879. Both stopped at step 500
without convergence, and optimizer continuation was unavailable. This
difference is not evidence for structural benefit. Moreover, H0 and H2
had identical encoder tensors at selected step-100 checkpoints because
the encoder warm-up lasted 1,600 steps. The interface H2 representation
was therefore an H0 alias, not an independent estimate of JEPA transfer.

The full-host engineering measurement, not a universal mechanism result,
gave 199,494,385,321 FLOPs for forward/backward profiling and a
shared-parameter gradient cosine of −0.052939 for one batch. Exact DP
and Python/C decoding were outside that profiler estimate.

\hypertarget{44-bpfold-motif-control}{%
\subsubsection{4.4 BPfold motif control}\label{44-bpfold-motif-control}}

Inner/outer motif input produced macro-F1 0.8514596293, identical to P0.
Shuffled motif reached 0.8515088595, exceeding inner/outer by
0.0000492302. Turner=0 gave 0.8476000946, and energy-only input equaled
P0. The decision was \texttt{STOP\_BPFOLD\_DEV\_NEGATIVE}; external
testing was not started. These results reject a benefit of this specific
input experiment, not the original BPfold method generally.

\hypertarget{45-interface-probes}{%
\subsubsection{4.5 Interface probes}\label{45-interface-probes}}

\textbf{Table 4. Frozen-representation interface results, mean ± SD over
three seeds}

\begin{longtable}[]{@{}lrr@{}}
\toprule
Representation/control & micro-AUPRC & Binary macro-F1 \\
\midrule
\endhead
RiNALMo & 0.7216 ± 0.0059 & 0.6448 ± 0.0076 \\
H0 host & 0.7715 ± 0.0053 & 0.6854 ± 0.0041 \\
Handcrafted local sequence features & 0.5438 ± 0.0070 & 0.4827 ±
0.0099 \\
Within-chain label shuffle & 0.6292 ± 0.0035 & 0.5429 ± 0.0025 \\
Length & 0.5522 ± 0.0132 & 0.3996 ± 0.0905 \\
GC content & 0.5524 ± 0.0465 & 0.3466 ± 0.0243 \\
\bottomrule
\end{longtable}

All 18 runs reached internal-validation plateaus, and independent metric
differences were at most 2.22×10⁻16. H0 outperformed the other
sequence-only controls on micro-AUPRC, but the within-chain shuffle
reached 0.6292, indicating information in inter-chain positive-rate or
composition differences. H0 AUPRC was 0.6233 on 100--499 nt chains
(n=23) versus 0.7725 on chains at least 1,500 nt long (n=88). These
results do not establish complex-environment understanding, independent
JEPA transfer, or improved secondary-structure prediction. Complete F2
did not pass because the structure-assisted arm was not computable with
the original full-pair ResNet2D.

\hypertarget{5-discussion}{%
\subsection{5. Discussion}\label{5-discussion}}

The audits demonstrate why downstream splitting and pretraining exposure
must be distinguished. Exact overlap is already sufficient to reject a
clean-OOD statement, while approximate-homology analysis exposes
additional matches that exact filtering cannot address. The family-level
zero for bpRNA-new and the 120 shared families for Rfam show that label
design and sequence exposure provide different evidence.

The model results also separate engineering success from scientific
support. The dual-view implementation passed mask, gradient, deployment,
and CUDA tests, yet H2 did not outperform its host or controls. The
motif implementation used verified official energy assets, yet produced
no improvement over P0 or shuffling. H0's interface result is a
qualified representation result limited by composition effects, length
strata, direct pretraining exposure, and missing structure-assisted
comparisons. None of these outcomes justifies reviving stopped
augmentation, RIPE, local Pair-JEPA, ensemble voting, or a 19,000 nt
beam to create a success narrative.

The mechanism gate was therefore narrowed rather than expanded. H0
already learns directly from pair labels, so using the same
canonical-nested targets in an additional latent objective does not, by
itself, establish a new structural mechanism. Pretraining-exposure
audits constrain ``unseen/OOD'' interpretation but do not block a
same-initialization comparison. No candidate with both an independent
information advantage and a falsifiable mechanism was identified, and no
new pilot was started. This is a scope-specific innovation decision, not
a universal claim that structural information is unhelpful.

\hypertarget{6-threats-to-validity}{%
\subsection{6. Threats to validity}\label{6-threats-to-validity}}

\begin{itemize}
\tightlist
\item
  \textbf{Incomplete historical corpus.} The NCBI nt URL has no
  historical release, date, or content hash. The official RiNALMo
  repository's \texttt{remote\_data.json} points to live resources;
  commit timestamps cannot identify the training snapshot. Current nt
  was not downloaded as a substitute. Complete historical exposure
  denominators remain unavailable.
\item
  \textbf{Search and coverage limits.} MMseqs2 is heuristic, and
  \texttt{-\/-max-seqs\ 300} is not a final output cap.
  \texttt{cov-mode=0} documentation alone was not accepted as proof of
  bidirectional coverage; independent postfiltering was applied. Short
  or low-complexity exact-positive sequences may not be recovered by
  MMseqs2, so exact and homology evidence are reported separately.
  No-hit is not proof of absence. Queries outside the 16--8192 nt filter
  are reported rather than counted as negatives.
\item
  \textbf{Family-label coverage.} Rfam 14.9 mapping covers only exact,
  uniquely mapped sequences. Ambiguity, missing metadata, later
  releases, and mixed TR1c sources prevent a complete
  family-independence claim.
\item
  \textbf{Protocol limits.} The dual-view pilot used one model-dev seed
  and a warm-start recipe. Native-initialization arms were not
  converged, and optimizer continuation was blocked. H2's long encoder
  warm-up prevents an independent interface JEPA estimate.
\item
  \textbf{Computational limits.} The original full-pair ResNet2D ran out
  of memory at 7,224 nt. Encoder-only long-sequence execution does not
  establish long-sequence accuracy or complete structure-feature
  support.
\item
  \textbf{Label scope.} Canonical-nested labels do not model complete
  three-dimensional structures or complex environments. TR1c mixes
  sources and is not an experimental-quality hierarchy. RNA--protein
  contacts are not RNA--RNA pairs, and unobserved residues were not used
  as negative labels.
\item
  \textbf{Downstream confounding.} Inter-chain composition can support
  nontrivial label-shuffle performance. Interface holdout had 86/202
  exact matches with RNAcentral, so downstream PDB/cluster exclusion
  alone cannot establish unseen-complex generalization.
\item
  \textbf{Calibration and pair dependence.} Pair-level ECE is not used
  as evidence of calibrated structure posteriors, and no structure-level
  posterior calibration is offered. Correlated pairs limit the
  interpretation of marginal calibration; binary interface F1 and AUPRC
  do not address that limitation.
\end{itemize}

\hypertarget{7-reproducibility-and-evidence-availability}{%
\subsection{7. Reproducibility and evidence
availability}\label{7-reproducibility-and-evidence-availability}}

The main code workspace is
\texttt{/home/rna\_jev\_xens2\_dualview\_round56}. Large evidence files
are under
\texttt{/teams/HoulabData\_1788502323/liucunyu/rinalmo\_pretrain\_audit}
and
\texttt{/teams/HoulabData\_1788502323/liucunyu/rna\_jev\_xens2\_dualview\_round56/artifacts/v2}.
The contract's Section 13 contains full commands and the rule that
previous result directories must be archived rather than overwritten.

Key artifacts include \texttt{exact\_overlap.json} (SHA256
\texttt{6cd88fc5\ldots{}5877}), \texttt{exact\_union\_summary.json}
(\texttt{d4412393\ldots{}baa4}), \texttt{homology\_union\_summary.json}
(\texttt{69a0b13b\ldots{}7131}),
\texttt{rfam14\_9\_exact\_sequence\_family\_mapping\_summary.json}
(\texttt{bdef1812\ldots{}8a99}),
\texttt{converged\_baseline\_summary.json}
(\texttt{2b465cd3\ldots{}4870}), \texttt{nt\_snapshot\_assessment.json}
(\texttt{5f4282bd\ldots{}6232}),
\texttt{xens2\_exposure\_stratified.json}
(\texttt{c1c39477\ldots{}07ec}), and
\texttt{historical\_bidirectional\_coverage.json}
(\texttt{8688a3fa\ldots{}62e5}). The final evidence presentation
contains 30 slides and has SHA256
\texttt{1b4b5432758290c93d4d95063b257f751ed5d84694bd4f6c79c6a004be832bbd}.

Relevant source revisions include \texttt{181b4d0} (tied-score AUPRC
correction), \texttt{77018f6} (H0/H2 encoder identity), \texttt{33805b9}
(bounded-memory homology processing), \texttt{610ee8e} (family mapping),
\texttt{9f586b6} (homology-union evidence page), \texttt{96aa207} (nt
provenance limitation), \texttt{e6dbffd} (frozen-prediction exposure
stratification), \texttt{8e00db0} (independent bidirectional-coverage
review), and \texttt{d23004c} (mechanism-novelty gate). Earlier
revisions are retained for model and engineering provenance. No
credentials or \texttt{.env} files are part of the artifacts.

\hypertarget{8-conclusion}{%
\subsection{8. Conclusion}\label{8-conclusion}}

Exact and approximate-homology analyses found extensive overlap between
the evaluated sequences and three public sources associated with RiNALMo
pretraining. The tested structure-assisted and motif variants did not
improve model-dev secondary-structure macro-F1, while the H0 interface
representation showed qualified transfer relative to sequence and
composition controls. These findings support contamination-aware
evaluation, explicit initialization and label audits, and clear
separation of engineering and scientific claims. They do not support a
new architecture accuracy breakthrough or an independent-generalization
claim. No additional pilot is justified under the current mechanism
gate. Future work must first establish an independent mathematical
mechanism with matched label scope, rather than expand data collection
or auditing to create an accuracy narrative.

\hypertarget{appendix-a-scope-specific-responses-to-anticipated-objections}{%
\subsection{Appendix A. Scope-specific responses to anticipated
objections}\label{appendix-a-scope-specific-responses-to-anticipated-objections}}

This appendix records landing points for the project's standing reviewer
questions. ``Not evaluated'' is a limitation, not evidence that the
corresponding claim holds.

\begin{longtable}[]{@{}ll@{}}
\toprule
\# & Landing point and scope-specific answer \\
\midrule
\endhead
\textbf{Q1} & Section 1 and Section 4.3. This is not a new CONTRAfold or
calibrated decision-head claim. The historical model lineage is
retained, and no System-1 calibration or matched-FLOPs advantage is
claimed. The dual-view comparison is limited to H0--H6 on model-dev; no
head-to-head evidence for a new decision paradigm is offered. \\
\textbf{Q2} & Sections 3.3 and 6. No exact-marginal calibration, RLCD
objective, or equality of pair probabilities with calibrated posteriors
is claimed. These were not evaluated in the present v2 protocol. \\
\textbf{Q3} & Section 6 contains a dedicated calibration-scope
limitation. Pair-level ECE is not used as evidence, and no
structure-level posterior calibration is offered. Pair correlation and
the absence of calibration evidence are explicit limitations. \\
\textbf{Q4} & Sections 3.3 and 4.4. The control is ``Nussinov + Turner
堆叠能'', not ViennaRNA. Turner=0 gives 0.8476000946, and no claim is
made that the baseline or motif input reproduces full ViennaRNA
thermodynamics. \\
\textbf{Q5} & Sections 3.4 and 6. The memory limitation is measured
rather than hidden: the original full-pair ResNet2D ran out of memory at
7,224 nt. No banded approximation, 19,000 nt beam, or full-length
effectiveness claim is used. \\
\textbf{Q6} & Sections 3.3 and 4.5. RiNALMo itself is the backbone and a
comparator, not a dismissed alternative. H0 and native RiNALMo are
evaluated separately, but H2 is not an independent JEPA estimate, and no
new backbone champion is claimed. \\
\textbf{Q7} & Sections 1 and 6. Canonical-nested pairs are the explicit
scope. General pseudoknots, noncanonical interactions, and full
three-dimensional or complex-environment modeling are not claimed, and
no H-type accuracy result is invented for v2. \\
\textbf{Q8} & Decision ordering is reported as an ablation; we make no
contribution claim about it. No new ordering experiment is presented in
v2. \\
\textbf{Q9} & Sections 3.1, 4.1, 4.2, and 6. Exact, 80\%-homology, and
official-family analyses are reported separately. Downstream cleaning
does not substitute for pretraining decontamination. No formal
raw/own-clean/matched-clean accuracy comparison is claimed because the
primary pilot did not pass. \\
\textbf{Q10} & Sections 2 and 3.3. Exact noncrossing DP is a preserved
baseline, not a new contribution. No illegal-structure-rate or
legalization-cost claim is made. \\
\textbf{Q11} & Section 4.3 and Section 7. No 50× claim, McCaskill
comparison, LinearPartition comparison, or new raw-sequence speed
benchmark is made. The 199,494,385,321-FLOP figure is a one-batch
profiler measurement, not a latency or speedup result. \\
\textbf{Q12} & Section 6 and the project's \texttt{limitations.md}
L1/L2. Jev's unreviewed status and undisclosed RLCD algorithm are not
used as evidence. No RLCD reproduction or vendor performance claim is
made; model claims are tied to the listed artifacts and official
software sources. \\
\bottomrule
\end{longtable}

\hypertarget{references}{%
\subsection{References}\label{references}}

{[}1{]} \emph{RiNALMo: RNA language model}. Nature Communications, 2025.
DOI: \texttt{10.1038/s41467-025-60872-5}.

{[}2{]} \emph{Mol-JEPA: A multimodal molecular world model based on
Joint Embedding Predictive Architecture}. arXiv:2608.22642v4. Official
repository commit \texttt{27f0065b69c73d18eec7a32ecc413b6dc3f33d90}; CC
BY-NC 4.0.

{[}3{]} \emph{BPfold project}. Nature Communications, 2025. Official
\texttt{key.energy}, MIT license; SHA256 prefix
\texttt{b657c553\ldots{}ce8b}.

{[}4{]} Danaee et al. \emph{bpRNA}. Nucleic Acids Research, 2018. DOI:
\texttt{10.1093/nar/gky285}.

{[}5{]} \emph{RNAcentral 22.0}. Official sequence archive:
\texttt{https://ftp.ebi.ac.uk/pub/databases/RNAcentral/releases/22.0/sequences/}.

{[}6{]} \emph{Rfam 14.9}. Official FASTA and database tables:
\texttt{https://ftp.ebi.ac.uk/pub/databases/Rfam/14.9/}.

{[}7{]} \emph{Ensembl release 109 and Ensembl Bacteria release 56}.
Official FTP archives:
\texttt{https://ftp.ensembl.org/pub/release-109/fasta/} and
\texttt{https://ftp.ensemblgenomes.ebi.ac.uk/pub/bacteria/release-56/fasta/}.

{[}8{]} \emph{MMseqs2}, version 17.b804f, installed from the Bioconda
channel. \texttt{https://github.com/soedinglab/MMseqs2}.

\end{document}
